\noindent
The general notion of minimality introduced in
Definition \ref{definition-minimal-versal}
can sometimes be deduced from the behaviour on tangent spaces.
Let $\xi$ be a formal object of the predeformation category
$\mathcal{F}$ and let
$\underline{\xi} : \underline{R}|_{\mathcal{C}_\Lambda} \to \mathcal{F}$
be the corresponding morphism. Then we can consider the following
the condition
\begin{equation}
\label{equation-bijective}
d\underline{\xi} : \text{Der}_\Lambda(R, k) \to T\mathcal{F} 
\text{ is bijective}
\end{equation}
and the condition
\begin{equation}
\label{equation-bijective-orbits}
d\underline{\xi} : \text{Der}_\Lambda(R, k) \to T\mathcal{F} 
\text{ is bijective on }\text{Der}_\Lambda(k, k)\text{-orbits.}
\end{equation}
Here we are using the identification
$T\underline{R}|_{\mathcal{C}_\Lambda} = \text{Der}_\Lambda(R, k)$ of
Example \ref{example-tangent-space-prorepresentable-functor}
and the action (\ref{equation-action}) of derivations
on the tangent spaces. If $k' \subset k$ is separable, then
$\text{Der}_\Lambda(k, k) = 0$ and the two conditions are equivalent.
It turns out that, in the presence of condition (S2) a versal formal
object is minimal if and only if $\underline{\xi}$ satisfies
(\ref{equation-bijective-orbits}).
Moreover, if $\underline{\xi}$ satisfies (\ref{equation-bijective}), then
$\mathcal{F}$ satisfies (S2).